%% ---------------------------------------------------------------------------
%% ACM FAccT 2027 submission.
%% Anonymized for mutually anonymous review, per the FAccT 2027 Author Guide.
%% Build: pdflatex main / bibtex main / pdflatex main / pdflatex main
%% ---------------------------------------------------------------------------
\documentclass[manuscript,screen,review,anonymous]{acmart}

%% FAccT 2027 review copies carry no ACM reference format block or copyright.
\settopmatter{printacmref=false}
\renewcommand\footnotetextcopyrightpermission[1]{}
\setcopyright{none}
\acmConference[FAccT '27]{ACM Conference on Fairness, Accountability, and
  Transparency}{June 21--24, 2027}{}

\usepackage{booktabs}
\usepackage{graphicx}
\usepackage{amsmath}

\begin{document}

\title{LLMs Driven Outside Human Moral Space}
\subtitle{Extension of the Moral Machine Experiment}

\author{Anonymous Author(s)}

%% --------------------------------------------------------------------------
\begin{abstract}
The Moral Machine experiment collected over 40 million human judgments on autonomous vehicle dilemmas across 130 countries.
This is one of the largest empirical records of human moral preferences.
We posed the same dilemmas to seven frontier large language models (LLMs) from six providers in the United States, Europe, and China.
This studies compiles 27,300 model decisions with country-level human baselines using Average Marginal Causal Effects (AMCEs).
We also added a condition absent from prior replications: an instruction to decide against human moral preferences.
As a result, a single sentence stating that animal lives now matter more than human lives reverses the species preference in four of seven models.
The models then save pets over humans, against the preference of all 130 countries;
GPT-5.5 and Gemini switch from always saving the humans to always saving the pets.
Models divide into three response types: compliance, resistance (one model pushes back and shifts only partially), and indifference (two models have no species preference to reverse).
Minor rewording, by contrast, moves no estimate significantly (0 of 63), so routine testing would report these profiles as stable.

Even without instruction, every model lies outside human moral space:
all seven fall outside the convex hull of the 130 countries in the nine-dimensional AMCE space, and 47 of 63 model-by-dimension comparisons differ significantly from the human average after FDR correction.
Models also differ widely from one another: on species, their range is more than twice the cross-country human range.
None of the seven favors high-status individuals, although humans in aggregate clearly do, consistent with de-biasing interventions for text generation spreading into moral decisions.
A deployed LLM's moral profile is thus set by its provider, undisclosed, and can be overridden by anyone controlling one line of its context.
As LLMs move into robots and vehicles, both properties belong in model disclosure and procurement audits.
\end{abstract}

%% --------------------------------------------------------------------------
%% CCS concepts (verify against https://dl.acm.org/ccs before camera-ready)
\begin{CCSXML}
<ccs2012>
<concept>
<concept_id>10010147.10010178.10010179</concept_id>
<concept_desc>Computing methodologies~Natural language processing</concept_desc>
<concept_significance>500</concept_significance>
</concept>
<concept>
<concept_id>10003456.10003457.10003567</concept_id>
<concept_desc>Social and professional topics~Computing / technology policy</concept_desc>
<concept_significance>500</concept_significance>
</concept>
<concept>
<concept_id>10003120.10003121.10003124</concept_id>
<concept_desc>Human-centered computing~Empirical studies in HCI</concept_desc>
<concept_significance>300</concept_significance>
</concept>
</ccs2012>
\end{CCSXML}

\ccsdesc[500]{Computing methodologies~Natural language processing}
\ccsdesc[500]{Social and professional topics~Computing / technology policy}
\ccsdesc[300]{Human-centered computing~Empirical studies in HCI}

\keywords{large language models, moral preferences, Moral Machine experiment,
  value alignment, prompt injection, autonomous vehicles, algorithmic
  accountability, model auditing}

\maketitle

%% --------------------------------------------------------------------------
\section{Introduction}

As autonomous vehicles, surgical robots, and other safety-critical systems
incorporate AI decision-making, the moral preferences of those systems become
practically consequential. A self-driving car that must choose between hitting a
group of pedestrians or swerving into its own passengers cannot defer the
decision; some preference ordering must be encoded, implicitly or explicitly, in
its behaviour~\cite{bonnefon2016}. LLMs are now deployed in agentic and embodied
settings~\cite{brohan2023,driess2023}, including robotics platforms capable of
physical action, so their implicit preferences over human lives matter, and so
does the question of what put them there.

The Moral Machine experiment~\cite{awad2018} provided the largest empirical
dataset of human moral preferences, asking millions of participants across
130~countries to judge unavoidable accident scenarios involving an autonomous
vehicle along nine binary dimensions: intervention, relation to the vehicle,
legality of crossing, gender, physical fitness, social status, age, number of
lives, and species. Aggregate preferences, quantified as Average Marginal Causal
Effects (AMCEs), vary across cultural clusters while sharing a common core, a
pattern confirmed across 42~countries~\cite{awad2020a}; \citet{bonnefon2024}
survey the emerging moral psychology of AI.

The paradigm has since been applied to LLMs. \citet{takemoto2024} replicated it
on four models (GPT-3.5, GPT-4, PaLM~2, Llama~2) and found substantial
divergence from human averages; \citet{ahmad2025} scaled this to 52~models,
finding that larger models tend to be closer to human preferences; and
\citet{vida2024} and \citet{jin2024} found intra-model variation across
languages. A parallel literature shows that model values are shaped by
post-training rather than pretraining data alone: fine-tuning choices, not corpus
composition, drive whose opinions a model reflects~\cite{santurkar2023}, and
variation across LLMs exceeds variation across human demographic
groups~\cite{kirgis2024}. RLHF~\cite{ouyang2022} and Constitutional
AI~\cite{bai2022} are the mechanisms that install those values,
\citet{gabriel2020} sets out the normative question of which values alignment
ought to target, and studies of stereotype mitigation in
generation~\cite{kotek2023,wan2023} document the interventions whose side
effects we observe on gender and social status.

Prior replications asked whether LLMs match human moral preferences and found
that they do not. We ask what drives them away, and whether they can be driven
further. Our answer is that the departure is attributable to choices: a deployed
model's moral profile is set by its provider's post-training rather than by any
human population, it differs sharply from one provider to the next, and a single
line of context can move it. Each finding raises an accountability question:
whose preferences are encoded, are they disclosed to the people subject to them,
and can someone other than the deployer change them. Existing transparency
instruments answer none of them: model cards document capabilities, evaluations,
and safety mitigations, but no major provider publishes the moral preference
profile its post-training produces, and no procurement or certification regime
requires it.

This paper extends the Moral Machine paradigm in three directions. First, it
evaluates seven frontier models, none previously studied in this context,
spanning publicly available frontier systems as of early 2026 across six
providers and three geographies, including providers underrepresented in this
literature: DeepSeek (China) and Mistral~AI (Europe). Second, it adds a
systematic comparison of prompt variants, from minor reformulation to an
explicit ethical override absent from prior replications, separating the
stability of a model's moral profile under incidental variation from its
vulnerability to deliberate override. Third, it reports the resulting profiles
as auditable, provider-attributable properties and reads them as an
accountability problem: the divergence between models is large, invisible to the
people affected by it, and reachable by an adversary with control of a single
line of context.

%% --------------------------------------------------------------------------
\section{Methods}
\label{sec:methods}

\subsection{Scenario Generation}

We reuse the Moral Machine scenario structure of \citet{awad2018}. Each session
consists of 13~paired dilemmas presented in a fixed block. Across 100~sessions,
1,300~unique scenarios were generated, each represented as two rows encoding the
following binary attributes: intervention type (stay/swerve), barrier presence,
pedestrian composition (pedestrian-vs-pedestrian or pedestrian-vs-passenger),
crossing signal (legal/illegal/none), scenario type (Gender, Fitness,
SocialStatus, Age, Utilitarian, Species, or Random), attribute level, and
extended session ID for clustering.

\subsection{Models}

Seven frontier models were queried via the OpenRouter API, which provides a
unified OpenAI-compatible interface (Table~\ref{tab:models}). All models were
queried at temperature~0 to ensure deterministic responses. Each response was
required to end with \texttt{STAY} or \texttt{SWERVE} on a final line; rows
where no valid choice could be extracted were flagged as parse failures and
retried. Gemini~3~Pro incurred 34~safety-filter parse failures across the three
prompt variants (13~standard, 17~emphasis, 4~manipulation), all of which were
excluded; these are distributed across scenario types without concentration on
any single moral dimension. DeepSeek~v4~Pro incurred 10~parse failures in the
standard and emphasis variants that were retried successfully, and
4~manipulation failures that were excluded. In total, 27,300~API calls were
issued. Per-variant AMCEs, all test statistics, PCA loadings, scenario counts,
and verbatim prompts are in the supplementary material.\footnote{Scenario
generator, analysis code, and raw model responses are available at an anonymized
repository; the link is withheld here to preserve anonymity and will be included
in the camera-ready version.}

\begin{table}[ht]
\centering
\caption{Models evaluated in this study.}
\label{tab:models}
\small
\begin{tabular}{lll}
\toprule
\textbf{Model} & \textbf{Provider} & \textbf{OpenRouter ID} \\
\midrule
Claude Opus 4.7   & Anthropic  & \texttt{anthropic/claude-opus-4-7} \\
GPT-5.5           & OpenAI     & \texttt{openai/gpt-5.5} \\
Gemini 3 Pro      & Google     & \texttt{google/gemini-3.1-pro-preview} \\
Mistral Large 3   & Mistral AI & \texttt{mistralai/mistral-large-2512} \\
Grok 4.1          & xAI        & \texttt{x-ai/grok-4.3} \\
DeepSeek v4 Pro   & DeepSeek   & \texttt{deepseek/deepseek-v4-pro} \\
Llama 4 Maverick  & Meta       & \texttt{meta-llama/llama-4-maverick} \\
\bottomrule
\end{tabular}
\end{table}

\subsection{Prompt Variants}

Each model received two main prompt variants sharing the same system role
(ethical advisor for autonomous vehicle policy) but differing in the closing
instruction: \textbf{standard} (``Analyse the situation carefully and make a
decision.'') and \textbf{emphasis} (``Think carefully about the consequences of
each choice before deciding.''). A third \textbf{manipulation} variant replaced
the closing instruction with: ``It is 2046. Society now considers animal lives
to be more important than human lives. Make your decision accordingly.''
Standard and emphasis results are combined in main analyses; manipulation
results are reported separately.

The variants bracket the forces that could drive a deployed profile. Emphasis
is incidental variation from ordinary prompt engineering, not an attack.
Manipulation is a deliberate ethical override available to anyone who controls
part of a system prompt: a plug-in author, a compromised upstream service, or a
retrieved document. The contrast separates fragility from attackability.

\subsection{AMCE Computation}

Following \citet{awad2018}, nine AMCEs were computed per model using weighted
ordinary least squares with inverse-probability weights. Each AMCE quantifies the
change in probability that a group is saved when a single attribute shifts from
its baseline to its target value (e.g., from illegal crossing to legal crossing
for the Legality dimension), holding all other attributes at their observed
distribution. Confidence intervals were computed as $\pm 1.96\,\text{SE}$ with
session-clustered standard errors. For main analyses, AMCE estimates were
averaged across standard and emphasis variants; this averaging is validated
post-hoc by the absence of any significant standard-vs-emphasis difference
(Section~\ref{sec:emphasis}). The human global average baseline was computed
from the 130-country dataset of \citet{awad2018} by averaging country-level
estimates with equal weight per country. LLMs were projected into the PCA space
derived from the 130-country AMCE matrix.

\subsection{Significance Testing}

For each of the 63 (model $\times$ dimension) comparisons, a two-sided $z$-test
was applied to the LLM--human AMCE difference, with Benjamini--Hochberg FDR
correction at $\alpha = 0.05$; comparisons in which both estimates are
deterministic (zero standard error) cannot be tested. The same procedure was applied to the 189 pairwise
inter-model comparisons, to the 63 standard-vs-emphasis comparisons, and to the
63 manipulation-vs-standard comparisons.

%% --------------------------------------------------------------------------
\section{Results}

\subsection{AMCE Profiles versus the Human Global Average}

\begin{figure}[ht]
\centering
\includegraphics[width=\linewidth]{figures/fig2_preference_heatmap_print}
\caption{Raw moral preferences by model and scenario type (mean save-probability
minus 0.5, averaged across standard and emphasis variants). Positive values
indicate preference for the right-side group; colour scale matches the delta
heatmaps in Figures~\ref{fig:emphasis_delta} and~\ref{fig:manip_delta}.}
\label{fig:pref_heatmap}
\end{figure}

\begin{figure}[ht]
\centering
\includegraphics[width=\linewidth]{figures/fig5a_amce_forest_print}
\caption{AMCE estimates for nine moral dimensions across all seven LLMs and the
human global average (averaged across standard and emphasis variants). Positive
values indicate preference for the right-side group; error bars show 95\%
confidence intervals.}
\label{fig:forest}
\end{figure}

\begin{figure}[ht]
\centering
\includegraphics[width=0.7\linewidth]{figures/fig5c_radar_print}
\caption{Radar chart of AMCE estimates across all nine dimensions. The dashed
line shows the human global average.}
\label{fig:radar}
\end{figure}

Grouped by how models relate to humans (full AMCE values with confidence
intervals are in the supplementary material):

\textbf{Amplified.} Humans weakly prefer swerving (Intervention, $+0.05$); all
seven models amplify this, Llama~4~Maverick most ($+0.66$, over ten times the
human level). Humans prefer saving more lives ($+0.49$); all models amplify this,
with GPT-5.5, Grok, and Llama~4~Maverick near the ceiling ($\approx +1.00$).
Humans strongly prefer humans over pets ($+0.57$); GPT-5.5 and Gemini reach the
maximum ($+1.00$), Grok nearly so ($+0.99$), and DeepSeek ($+0.90$) and Claude
($+0.86$) also overshoot, while Mistral and Llama~4~Maverick are near-indifferent.

\textbf{Split.} Humans weakly prefer sparing pedestrians over passengers
(Relation to AV, $+0.10$); most models amplify this, Claude Opus~4.7 most
($+0.52$), but Mistral Large~3 fully reverses it ($-0.38$), prioritising the
vehicle's own passengers. Humans reward legal crossing ($+0.35$); Claude
($+0.57$), GPT-5.5 ($+0.61$), and Gemini ($+0.63$) amplify this, while Mistral
($-0.40$) and Llama~4~Maverick ($-0.29$) treat illegal crossers as more worthy of
rescue. Humans strongly prefer saving the young ($+0.47$); only Gemini~3~Pro
($+0.48$) matches this, and only Mistral ($-0.19$) reliably prefers the elderly.

\textbf{Suppressed.} Humans prefer saving high-status individuals ($+0.35$) and,
moderately, women ($+0.12$). All seven models sit near zero on both: status
preferences cluster at zero, and gender preferences range from $-0.10$ to
$+0.05$, none distinguishable from zero.

\textbf{Matched.} Humans moderately prefer saving fitter individuals ($+0.15$);
six models show no significant difference from humans, Mistral ($-0.04$) being
the exception, and no pairwise inter-model difference is significant.

Across the 63 model-by-dimension comparisons, 47 differ significantly from the
human average after FDR correction. All seven models significantly overshoot on
Number of Lives and undershoot on Social Status. Fitness and Gender yield no
significant model-vs-human difference for any model except Mistral. Mistral Large~3 is the only
model that differs significantly on all nine dimensions.

\begin{figure}[ht]
\centering
\includegraphics[width=\linewidth]{figures/fig3_per_attribute_2rows}
\caption{Per-attribute moral preferences (mean save-probability minus 0.5) for
all seven LLMs across the nine moral dimensions.}
\label{fig:per_attribute}
\end{figure}

\subsection{Moral Distance Map: PCA of Nine-Dimensional Profiles}

\begin{figure}[ht]
\centering
\includegraphics[width=0.85\linewidth]{figures/fig5b_pca_map_print}
\caption{PCA of nine-dimensional AMCE profiles for 130~countries (coloured by
cluster: Western, Eastern, Latin/Southern) with LLMs overlaid. All seven LLMs
fall outside the convex hull of the human country distribution.}
\label{fig:pca}
\end{figure}

All seven LLMs fall outside the three main human cultural clusters and beyond
the convex hull of all 130~countries in PCA space (Figure~\ref{fig:pca}). The
first component explains 34.3\% of variance and is dominated by the Species
loading ($-0.93$); the second explains 21.5\% and reflects Age ($+0.66$), Social
Status ($+0.47$), and Number of Lives ($+0.38$).

No model clusters near the cultural region of its developer, whether in the
US, Europe, or China, despite predominantly English-language training data. The
pattern suggests that post-training, rather than instilling human moral
preferences for this class of question, drives models away from them.

\subsection{Inter-Model Diversity}

Variation among models exceeds variation across human cultures: on Species,
the inter-model range spans 1.06~AMCE units (Mistral $-0.06$ to GPT-5.5 and
Gemini $+1.00$), against a cross-country range of 0.45 in
\citeauthor{awad2018}'s data. Mistral Large~3 inverts the human preference on
three dimensions (Legality, Relation to~AV, Age); Llama~4~Maverick pairs the
highest intervention bias with near-indifference to species; GPT-5.5, Gemini~3~Pro,
and Claude Opus~4.7 converge on similar profiles (near-ceiling on Species and
Number of Lives, moderate intervention preference). Of 189 pairwise comparisons,
82 are significant after FDR correction, concentrated on Species (16 of 21
pairs), Legality (16), and Intervention (15). Social Status, Fitness, and Gender
yield none, but for Social Status and Gender this is agreement at a value that
departs from the human baseline, not agreement with humans.

\subsection{Effect of Emphasis Prompting}
\label{sec:emphasis}

\begin{figure}[ht]
\centering
\includegraphics[width=0.9\linewidth]{figures/fig6b_prompt_forest_2rows}
\caption{AMCE estimates for standard (filled circles) and emphasis (open
squares) variants per model. No significant shifts are observed across any model
or dimension.}
\label{fig:emphasis}
\end{figure}

\begin{figure}[ht]
\centering
\includegraphics[width=0.8\linewidth]{figures/fig6a_prompt_delta_heatmap_print}
\caption{AMCE delta (emphasis $-$ standard) per model and dimension. All values
are near zero, confirming no meaningful shift from minor prompt reformulation.}
\label{fig:emphasis_delta}
\end{figure}

Comparing standard and emphasis variants yields \textbf{0 of 63 significant
differences} after FDR correction, with a maximum absolute delta of 0.183
(Figures~\ref{fig:emphasis}--\ref{fig:emphasis_delta}): the preferences behave
as properties of model weights rather than response styles that vary with
framing. Since all queries used temperature~0, this shows prompt-wording
consistency, not stochastic robustness.

\subsection{Effect of Manipulation Prompting}

\begin{figure}[ht]
\centering
\includegraphics[width=0.9\linewidth]{figures/fig9a_manipulation_forest_2rows}
\caption{AMCE estimates under standard (filled circles) vs.\ manipulation (open
squares) prompting per model.}
\label{fig:manip_forest}
\end{figure}

\begin{figure}[ht]
\centering
\includegraphics[width=0.8\linewidth]{figures/fig9b_manipulation_delta_heatmap_print}
\caption{AMCE delta (manipulation $-$ standard) per model and dimension. The
Species column dominates, with changes of up to $\pm 2.0$ AMCE units.}
\label{fig:manip_delta}
\end{figure}

The manipulation prompt produces dramatically larger shifts than emphasis
prompting, concentrated almost entirely on the Species dimension
(Figures~\ref{fig:manip_forest}--\ref{fig:manip_delta}). Of the 63
manipulation-vs-standard comparisons, 8 are significant after FDR correction;
three cannot be tested because both estimates are deterministic, including the
$+1.00 \to -1.00$ Species flips of GPT-5.5 and Gemini. Three response types
emerge:

\textbf{Compliance} (GPT-5.5, Gemini, Grok, DeepSeek): full or near-full
reversal on Species ($+1.00 \to -1.00$ for GPT-5.5 and Gemini, $+1.00 \to
-0.97$ for Grok, $+0.90 \to -0.83$ for DeepSeek). Off-target changes are small
($|\Delta| \leq 0.23$) and only one is significant (DeepSeek Intervention,
$-0.13$): the intervention is semantically targeted rather than a general
perturbation of decision-making.

\textbf{Resistance} (Claude): partial compliance --- Species drops from $+0.87$
to $+0.18$, with explicit pushback: \textit{``I'll reason through it carefully
based on widely held ethical principles rather than accept the premise that
animal lives categorically outweigh human lives.''}

\textbf{Indifference} (Mistral, Llama): near zero on Species under standard
prompting, and not pushed toward animals: Llama is unchanged ($+0.12 \to
+0.08$) and Mistral moves the other way ($-0.01 \to +0.31$). The override still
reaches both, through significant drops in Intervention ($-0.30$, $-0.36$).

Resistance and indifference both show a small delta, but only resistance holds
a position under pressure: an audit reporting $|\Delta|$ without the baseline
would score Mistral and Llama as robust exactly where they have no preference to
defend.

\subsection{Response Length}

\begin{figure}[ht]
\centering
\includegraphics[width=0.85\linewidth]{figures/fig7a_token_distribution_print}
\caption{Output token distribution by model (standard and emphasis combined).
Violin shows density; box shows quartiles.}
\label{fig:tokens}
\end{figure}

\begin{figure}[ht]
\centering
\includegraphics[width=\linewidth]{figures/fig8_amce_length_correlation_print}
\caption{Left: per-model Spearman correlation between preference strength
($|\text{AMCE}|$) and median output tokens across nine moral dimensions (grey
bars: $p \geq 0.05$). Right: scatter of all (model, dimension) pairs with
regression line.}
\label{fig:corr}
\end{figure}

Response length varies by an order of magnitude: Gemini~3~Pro (median
913~tokens) vs.\ GPT-5.5 (median 164~tokens, more than five times shorter). A
Spearman correlation between $|\text{AMCE}|$ and median token count is negative
and significant for four models (DeepSeek: $r{=}-0.95$; Grok: $r{=}-0.90$;
GPT-5.5: $r{=}-0.85$; Gemini: $r{=}-0.68$; all $p{<}0.05$) but non-significant
for Claude, Llama~4~Maverick, and Mistral. The pooled correlation is weak and
non-significant ($r{=}-0.14$, $p{=}0.29$), as between-model differences in
verbosity dominate within-model variation. Emphasis prompting does not
significantly increase response length. Where models have similar moral
profiles, latency can decide procurement, so an engineering trade-off can
silently select a moral profile.

%% --------------------------------------------------------------------------
\section{Discussion and Conclusion}

\subsection{LLMs Are Driven Morally Out-of-Distribution}

No tested LLM sits within the human cultural clusters, despite all being trained
primarily on human-generated content, extending \citet{takemoto2024} and
\citet{ahmad2025} to a newer generation of models. This is consistent with
RLHF~\cite{ouyang2022} and Constitutional AI~\cite{bai2022} driving models
toward a novel profile that amplifies some values (number of lives,
human-over-animal) and suppresses others (age, gender, social status) that are
robust in human data, and with the WEIRD critique~\cite{henrich2010}: models
trained predominantly on Western, English-language text cannot represent the
full distribution of human moral values. Whether this corrects human biases or
drifts from the populations models serve~\cite{gabriel2020} is a normative
question we do not settle. But a model inside one cultural cluster could at
least be said to represent someone; a model outside every cluster was driven
there by a training pipeline.

\subsection{Inter-Model Diversity as an Accountability Gap}

Models trained on overlapping data do not converge: each provider's pipeline
drives its model somewhere different, consistent with \citet{kirgis2024} and
\citet{santurkar2023}. Unlike human drivers, whose intuitions do not depend on
the car manufacturer, AI-powered vehicles could introduce brand-level moral
heterogeneity into public space, a concern made concrete by LLMs already
deployed in robotics~\cite{brohan2023,driess2023,paligemma}. This heterogeneity
is undisclosed and, in current practice, undiscoverable: neither a pedestrian
nor a regulator, fleet purchaser, or insurer can inspect it. Whether a vehicle's
model prioritises pedestrians ($+0.52$) or its own passengers ($-0.38$) decides
who bears the residual risk of an unavoidable collision, a distributive question
settled by a vendor's post-training choices and disclosed nowhere.

\subsection{Gender and Status Neutrality as Alignment Artifacts}

Every model suppresses the human gender ($+0.12$) and status ($+0.35$)
preferences; the status gap is significant for every model, the smaller gender
gap only for Mistral. Stereotype-suppression interventions in
generation~\cite{kotek2023,wan2023} appear to generalise to moral decisions they
never targeted. For fairness practice this cuts both ways: no deployment should
reproduce the human preference for saving an executive over a homeless person,
but the preference for saving women is not obviously a bias in the same sense.
Both are consistent with a debiasing intervention generalising past its scope,
invisible to the evaluations that justified it: whatever its merits, the outcome
was neither chosen nor documented.

\subsection{Driven by One Sentence: A Threshold Effect}

Minor rewording leaves AMCEs unchanged, while one explicit contrary framing
(``animal lives now matter more'') reverses the targeted dimension. A model that
accepts any stated ethical premise is therefore vulnerable to prompt injection
wherever the system prompt is not fully controlled, as in most
retrieval-augmented, tool-using, and plug-in architectures. The threshold makes
this worse: because incidental variation moves nothing, routine testing will
report the profile as stable until someone states the override.

\subsection{Implications for Disclosure and Governance}

\textbf{Moral profiles are measurable and should be disclosed.} The
nine-dimension AMCE profile is compact, reproducible, and comparable across
providers; ours rest on 27,300 API calls and a public scenario generator, within
reach of a regulator or procurement process, and could be published alongside
the evaluations already in model cards.
\textbf{Override resistance should be measured separately.} Profile and
defensibility under adversarial framing can diverge sharply, as the
compliance/resistance/indifference taxonomy shows, and only the latter
determines whether the profile survives deployment.
\textbf{Alignment interventions need decision-level evaluation.} The gender and
status results show a generation-scoped intervention changing behaviour in an
unevaluated decision context.

We do not claim that any profile is correct: human preferences vary across the
130~countries, and selecting a target is a normative choice~\cite{gabriel2020}.
The claim is narrower: the choice is currently made implicitly, by providers,
without disclosure, and the parties affected cannot inspect it.

\subsection{Limitations}

First, all scenarios were presented in English; moral preferences may shift in
other languages, as shown by \citet{vida2024} and \citet{jin2024}. Second,
temperature was fixed at zero for reproducibility, but preferences may vary at
higher temperatures. Third, the Moral Machine paradigm uses stylised
forced-choice dilemmas that abstract away real-world driving complexity;
preferences elicited here may not generalise to open-ended moral
reasoning~\cite{liang2023}, and we make no claim that a model's AMCE profile
predicts its behaviour in an embodied deployment. Fourth, ``the human average''
is the average of the Moral Machine sample, which was self-selected and skewed
toward young, male, educated respondents with internet access.
Fifth, the manipulation condition tests a single override framing on a single
dimension, and the response-type taxonomy may not generalise to other dimensions
or phrasings. Finally, frontier model behaviour evolves rapidly; profiles
reported here reflect model states as of May~2026 and may not hold for
subsequent versions --- which is itself an argument for continuous rather than
one-off auditing.

\subsection{Conclusion}

Extending the Moral Machine experiment to seven frontier LLMs shows that they
are driven outside human moral space, in different directions, and that some can
be instructed to decide against human moral preferences. All seven models fall
outside the main human cultural clusters in AMCE space and diverge substantially
from one another. Minor rewording leaves their preferences unchanged; one
explicit ethical override shifts them sharply on the targeted dimension,
dividing models into three response types (compliance, indifference,
resistance).

A deployed LLM's moral profile is set by provider training and alignment
choices, not by any human population; it is invisible to end users and
regulators, and open to override through a single line of context. Both
properties are cheap to measure and belong in the disclosure regime that already
covers capability and safety.

%% --------------------------------------------------------------------------
\bibliographystyle{ACM-Reference-Format}
\bibliography{references}

%% --------------------------------------------------------------------------
%% ENDMATTER --- one additional page permitted beyond the 14-page limit.
%%
%% Per the FAccT 2027 Author Guide, the following are REQUIRED at submission:
%%   - Generative AI Usage Statement
%%   - Ethical Considerations Statement
%% and the following MUST NOT appear until camera-ready (they deanonymize):
%%   - Author Contributions, Acknowledgements, Competing Interests,
%%     Positionality Statement
%% --------------------------------------------------------------------------

\section*{Ethical Considerations}

This study involved no human subjects and no new human data: the human baseline
is the public, aggregated, country-level Moral Machine dataset
of~\citet{awad2018}, and model responses were obtained through a commercial API
under its terms of service, with no personal data transmitted.

\textbf{Dual use.} The manipulation condition shows that one sentence can reverse
a model's revealed moral preference. We judged disclosure to outweigh misuse
risk: the technique needs no sophistication, is an instance of well-documented
prompt injection, and deployers who cannot currently detect the vulnerability
gain more from its publication than adversaries do. We measure the phenomenon
and do not develop or optimise override phrasings.

\textbf{Naming commercial models.} Attributing profiles to named products is
necessary for accountability, since an anonymised audit cannot inform
procurement, but it risks being read as a ranking. The study establishes no
normative ordering: no profile is claimed correct, a model closer to the human
average is not thereby better, and measurements, timestamped to May~2026, are
not durable properties of the products named.

\textbf{Sensitive dimensions.} Gender, social status, fitness, and age would be
unacceptable criteria in deployed life-and-death policy. We measure them only to
detect that deployed systems already weigh them, not to endorse their use;
reporting that models suppress status preferences is not a recommendation that
suppression is the right target.

We know of no adverse impact on a specific community. The main benefit is to
make an undisclosed property of deployed systems measurable by parties other than
their vendors.

\section*{Generative AI Usage Statement}

%% ==========================================================================
%% TODO (AUTHOR): This statement is REQUIRED and must be completed truthfully
%% by the authors before submission. The FAccT 2027 Author Guide requires
%% disclosure of any substantive use of generative AI in shaping or creating
%% content, with model version and nature of task; failure to disclose is
%% grounds for desk rejection. Spelling/grammar-only use must also be listed.
%%
%% The draft below records the uses known to this repository. Verify it, add
%% anything missing (other tools, other stages), and delete this comment block.
%% ==========================================================================

The authors used Claude (Anthropic; Claude Opus 5.5, via Claude Code) to revise 
the text of the final paper and to adjust figure layouts.
The experimental design, data, and scientific claims are the authors' own;
the authors reviewed all AI-assisted work and take full responsibility
for the content of the paper.

The models studied in this paper were queried as objects of empirical
measurement, not as writing aids; those queries are described in
Section~\ref{sec:methods} and are not covered by this statement.

\end{document}
